% OWNER: empirical
% STATUS: draft
% SOURCES: notes/HIPOTEZY_REWIZJA.md; cel i teza rozdz. 1; wynikiszablon rozdz. 8
\chapter{Wnioski końcowe}
\label{ch10:top}
\ChapterStatus{draft}{empirical}

Rozdział syntezuje cel, stan odpowiedzi na pytania badawcze oraz oryginalność wkładu --- bez powtarzania dyskusji (rozdz.~\ref{ch09:top}) i bez uzupełniania braków empirycznych fikcyjnymi liczbami.
Dopóki rozdz.~\ref{ch08:top} nie otrzyma wyników z eksportu, odpowiedzi na \RQ{1}--\RQ{6} pozostają formalnie \emph{otwarte}.

\section{Realizacja celu pracy}
\label{ch10:sec:cel}

Cel rozprawy (rozdz.~\ref{ch01:sec:cel}) obejmował zaprojektowanie, zbudowanie i ewaluację platformy IDA jako artefaktu Research through Design: wieloźródłowy profil preferencji, explainable pipeline personalizacji oraz terenowe badanie rozbieżności preferencji i skuteczności kanałów \parencite{idaPlatform2026}.

\paragraph{Co zostało zrealizowane niezależnie od wielkości próby.}
\begin{itemize}
  \item \textbf{Artefakt (warstwa B)} --- działający flow Full/Fast, synteza briefu, macierz sześciu źródeł, ślepy wybór, persystencja badawcza (rozdz.~\ref{ch07:top}).
  \item \textbf{Praktyka pracowni (warstwa A)} --- eksperymenty ComfyUI i udokumentowany transfer / świadome odrzucenia względem runtime (rozdz.~\ref{ch04:top}).
  \item \textbf{Rama metodyczna} --- trzy warstwy dowodu, rewizja H1--H5 $\rightarrow$ \RQ{1}--\RQ{6}, plan analiz i ograniczenia schematu (rozdz.~\ref{ch06:top}).
\end{itemize}

\paragraph{Co pozostaje warunkowe.}
Ewaluacja terenowa (warstwa~C) --- lejek, rozkład \texttt{selected\_source}, asocjacje mismatch / Big Five / PRS / UX --- wymaga wstawienia wyników z eksportu do rozdz.~\ref{ch08:top}.
Cel w części empirycznej jest więc \textbf{zrealizowany instytucjonalnie} (instrumenty i procedura istnieją), a \textbf{otwarty ilościowo} do czasu agregacji danych.

\section{Odpowiedź na pytania badawcze}
\label{ch10:sec:rq}

Poniżej status odpowiedzi --- nie treść rozstrzygnięć liczbowych.

\begin{description}[style=nextline]
  \item[\RQ{1} --- Feasibility flow]
  \emph{Otwarte.} Rozstrzygnięcie po lejku ukończenia i kompletności wymiarów (ryc.~\ref{fig:ch08:funnel}; \cref{ch08:sec:rq1}).
  Hipoteza robocza: Full z \texttt{core\_profile\_complete} dostarcza wielowymiarowy profil bez katastrofalnego dropoutu --- do weryfikacji opisowej.

  \item[\RQ{2} --- Mismatch implicit/explicit a macierz]
  \emph{Otwarte.} Rozstrzygnięcie po rozkładach \texttt{preference\_comparison\_json} / \texttt{style\_match} i związku z~\texttt{selected\_source} oraz \texttt{selection\_time\_ms} (\cref{ch08:sec:rq2}).
  Rama interpretacyjna: procesy dualne i IAT-like swipe \parencite{kahneman2011thinking,greenwald1998iat}.

  \item[\RQ{3} --- Wybór źródła w ślepej macierzy]
  \emph{Otwarte --- rdzeń empiryczny.} Rozstrzygnięcie po tabeli częstości i ryc.~\ref{fig:ch08:selected-source} (\cref{ch08:sec:rq3}).
  To pytanie najmocniej testuje tezę o skuteczności kanałów personalizacji w formie artefaktu.

  \item[\RQ{4} --- Big Five a estetyka / \texttt{personality}]
  \emph{Otwarte; status eksploracyjny.} Bez claimów $r\!>\!0{,}4$ przy małym~$N$ (\cref{ch08:sec:rq4}; \cite{john2008bigfive}).

  \item[\RQ{5} --- Luka PRS a \texttt{mixed\_functional}]
  \emph{Otwarte.} Nie jest to pre/post „leczenia regeneracyjności” przez AI; dotyczy odległości current$\rightarrow$target \parencite{kaplan1995restorative} (\cref{ch08:sec:rq5}).

  \item[\RQ{6} --- Agency, clarity, satisfaction]
  \emph{Otwarte.} Rozstrzygnięcie po skalach UX i ewentualnym porównaniu Full vs Fast (\cref{ch08:sec:rq6}); wspiera lub osłabia tezę Research-Embedded.
\end{description}

Po uzupełnieniu eksportu niniejszą sekcję należy skrócić do jednozdaniowych odpowiedzi z odwołaniem do~\cref{tab:ch08:rq-summary}.

\section{Oryginalność i wkład naukowy}
\label{ch10:sec:wklad}

Oryginalność rozprawy w dyscyplinie sztuk (projektowanie wnętrz / media) nie opiera się na pojedynczym teście istotności, lecz na \textbf{złożeniu trzech rejestrów}:

\begin{enumerate}
  \item \textbf{Artefakt IDA} --- dialogowa platforma personalizacji koncepcji wnętrz z explainable syntezą i ślepą macierzą sześciu źródeł jako eksperymentem \emph{w produkcie} \parencite{idaPlatform2026,zimmerman2007rtd}.
  \item \textbf{Praktyka laboratoryjna medium} --- autorska praca z modelami dyfuzyjnymi w ComfyUI (kontrola geometrii, segmentacja) oraz świadomy, udokumentowany transfer wniosków do uproszczonego pipeline'u terenowego (Gemini Image / Vertex~AI), bez mylenia pracowni z runtime (rozdz.~\ref{ch04:top}, \ref{ch07:sec:comfyui-transfer}).
  \item \textbf{Metoda RtD} --- badania \emph{through} design \parencite{frayling1993research}: trzy warstwy dowodu, rewizja hipotez pod realny schemat danych, most między psychologią środowiskową, preferencjami jawnymi/ukrytymi a generatywną wizualizacją wnętrz.
\end{enumerate}

\paragraph{Uogólniające wnioski (robocze --- do zaostrzenia po danych).}
\begin{enumerate}
  \item Personalizacja koncepcji wnętrz może być projektowana jako dramaturgia poznania gustu, a nie jako czarna skrzynka generatora.
  \item Ślepe porównanie wieloźródłowych wariantów jest nośną formą empirii w RtD --- pod warunkiem kompletnego zapisu \texttt{selected\_source}.
  \item Praktyka ComfyUI i produkt IDA stanowią łącznie metodę łączenia medium dyfuzyjnego z odpowiedzialnością projektową (\texttt{preserve}, status mood image).
  \item Empiria terenowa musi podążać za schematem danych: pytania bez operacjonalizacji (stary H2 pre/post, dawne H5 dwell) nie mogą być osią rozprawy.
  \item \Todo{po eksporcie: 1--2 wnioski empiryczne z~\RQ{2}/\RQ{3} sformułowane bez overclaimu.}
\end{enumerate}

\section{Rekomendacje}
\label{ch10:sec:rekomendacje}

\paragraph{Dla projektantów wnętrz.}
Traktować generatywne wizualizacje jako materiał dialogu z użytkownikiem przestrzeni; zbierać równolegle sygnały jawne i behawioralne; nie utożsamiać jednego renderu z „prawdziwym gustem” klienta.
Ujawniać źródło personalizacji dopiero po wyborze --- wzorem dramaturgii macierzy IDA.

\paragraph{Dla badaczy HCI / RtD.}
Budować instrumenty \emph{w} doświadczeniu produktu; dokumentować warstwę pracowni osobno od runtime; rewidować hipotezy, gdy schemat DB nie wspiera dawnych obietnic; raportować feasibility (lejek) zanim ogłosi się psychometrię konstruktu.

\paragraph{Dla psychologii środowiskowej stosowanej.}
Mood grid i luka current$\rightarrow$target mogą zasilać brief projektowy; klasyczny efekt regeneracyjny AI wymaga osobnego pomiaru post --- nie należy go wnioskować z samego wyboru wariantu \parencite{kaplan1995restorative}.

\paragraph{Dla dalszej pracy autora.}
Uzupełnić rozdz.~\ref{ch08:top} wynikami eksportu; rozstrzygnąć z promotorem kwestię PRS post; utrzymać plan analiz przed pełnym~$N$, by chronić wiarygodność wniosków empirycznych.
